\section{System Overview}
EQUITAS-RCMF takes two inputs: a $224 \times 224$ RGB facial video frame $X_V$ and a physiological feature vector $X_P$ derived from wrist-worn blood volume pulse (BVP) and electrodermal activity (EDA) sensors. The forward pass proceeds as
\begin{equation}
X_V \xrightarrow{\text{MobileNetV3}} \mathbf{v}_{\text{raw}}
\xrightarrow{\text{Stiefel}} (\mathbf{v}_{\text{causal}}, \mathbf{v}_{\text{conf}})
\xrightarrow{\text{Gate}} \mathbf{z}_{\text{fused}}
\xrightarrow{\text{Classifier}} \hat{Y},
\end{equation}
with the physiological embedding joining at the gating stage. The stages are described below.

\section{Vision Encoder}
The vision branch uses a MobileNetV3-Small backbone~\cite{howard2019} pretrained on ImageNet-1K. It produces a spatial feature map $\mathbf{A} \in \mathbb{R}^{576 \times h \times w}$, which adaptive average pooling reduces to a vector:
\begin{equation}
\mathbf{v}_{\text{raw}} = \operatorname{Pool}(\mathbf{A}) \in \mathbb{R}^{576}.
\end{equation}
% TODO(authors): The markdown states 1,140,133 trainable parameters in one place and
% "approximately 2.5M parameters" for MobileNetV3-Small in another. Report the full
% model's parameter count and, if part of the backbone is frozen, the trainable count,
% so the two figures are not read as contradictory.

\section{Physiological Encoder}
The physiological branch processes the BVP and EDA feature vector with a two-layer multilayer perceptron using layer normalization and SiLU activations, producing an embedding $\mathbf{p}_{\text{emb}} \in \mathbb{R}^{d}$ with $d = 192$.

\section{Stiefel Orthogonal Subspace Decomposition}
A learnable matrix $W_{\text{raw}} \in \mathbb{R}^{256 \times 576}$ is projected at every forward pass onto the set of matrices with orthonormal rows. Using the thin QR decomposition $W_{\text{raw}}^{\top} = QR$ with $Q \in \mathbb{R}^{576 \times 256}$~\cite{edelman1998},
\begin{equation}
W_{\text{St}} = Q^{\top}, \qquad W_{\text{St}} W_{\text{St}}^{\top} = I_{256}.
\end{equation}
The rows of $W_{\text{St}}$ are split into a causal block and a confounder block,
\begin{equation}
W_{\text{causal}} = W_{\text{St}}[1{:}192, :], \qquad W_{\text{conf}} = W_{\text{St}}[193{:}256, :],
\end{equation}
which gives $\mathbf{v}_{\text{causal}} = W_{\text{causal}}\mathbf{v}_{\text{raw}} \in \mathbb{R}^{192}$ and $\mathbf{v}_{\text{conf}} = W_{\text{conf}}\mathbf{v}_{\text{raw}} \in \mathbb{R}^{64}$. Because the rows of $W_{\text{St}}$ are orthonormal,
\begin{equation}
W_{\text{causal}} W_{\text{conf}}^{\top} = 0,
\end{equation}
so the two blocks project $\mathbf{v}_{\text{raw}}$ onto orthogonal subspaces. The orthogonality error $\|W_{\text{causal}} W_{\text{conf}}^{\top}\|_F$ was measured as $1.7149 \times 10^{-6}$ at initialization.

Orthogonality of the projections does not by itself make the two representations statistically independent: information about the scar can still be present in $\mathbf{v}_{\text{causal}}$, and both blocks are derived from the same parameters $W_{\text{raw}}$, so gradients from all loss terms update both~\cite{sarhan2020, huang2018}. The decomposition therefore provides a structured place for scar information to be routed, while the invariance losses in Section~\ref{sec:loss} discourage it from remaining in the causal subspace.

\section{LUPI Confounder Focus: Training and Inference Modes}
The model estimates how strongly the visual features concentrate on the scar region. Following the learning-using-privileged-information paradigm~\cite{vapnik2009, vapnik2015}, this estimate uses the ground-truth scar mask only during training.

\paragraph{Privileged mode (training).} With the scar mask $M$ available, the focus is the log-normalized ratio of feature-map energy inside the mask to the energy over the whole map:
\begin{equation}
\text{Focus} = \log\!\left(1 + \frac{\text{mean activation energy of } \mathbf{A} \text{ inside } M}{\text{mean activation energy of } \mathbf{A} + \epsilon}\right).
\end{equation}

\paragraph{Autonomous mode (inference).} No mask is provided. The focus is instead estimated from the confounder subspace:
\begin{equation}
\text{Focus}_{\text{auto}} = \sigma\bigl(\mathbf{w}_{\text{conf}}^{\top}\mathbf{v}_{\text{conf}} + b_{\text{conf}}\bigr).
\end{equation}
When the model is compiled for deployment in this mode, the mask input and the computations that depend on it are absent from the traced graph (Section~\ref{sec:deployment}).
% TODO(authors): State how Focus_auto is trained to match the privileged Focus (for
% example, a regression or distillation loss between the two), since the loss
% functional below does not include such a term explicitly.

\section{Focus-Driven Attention Gate}
The contribution of visual features is controlled by a gate that combines a learned cross-modal term with an exponential penalty on the focus:
\begin{equation}
G = \sigma\!\bigl(\operatorname{MLP}_{\text{cross}}([\mathbf{v}_{\text{causal}}; \mathbf{p}_{\text{emb}}])\bigr) \cdot \exp(-\kappa \cdot \text{Focus}),
\end{equation}
where $\kappa$ is a learnable temperature initialized at $1.5$ and clamped to $[0.1, 10.0]$. The factor $\exp(-\kappa \cdot \text{Focus})$ reduces the gate, and hence the weight given to visual features, when those features concentrate on the scar region.

\section{Latent Fusion and Classification}
The fused representation weights the visual and physiological embeddings by the gate,
\begin{equation}
\mathbf{z}_{\text{fused}} = G \odot \mathbf{v}_{\text{causal}} + (1 - G) \odot \mathbf{p}_{\text{emb}},
\end{equation}
and the classifier is
\begin{equation}
\hat{Y} = \operatorname{Softmax}\!\Bigl(\operatorname{Linear}_{128 \to 2}\bigl(\operatorname{Dropout}_{0.15}(\operatorname{SiLU}(\operatorname{Linear}_{192 \to 128}(\mathbf{z}_{\text{fused}})))\bigr)\Bigr).
\end{equation}

\section{Training Objective}
\label{sec:loss}
The model is trained with the combined objective
\begin{equation}
\mathcal{L}_{\text{total}} = \mathcal{L}_{\text{task}} + 0.5\,\mathcal{L}_{\text{conf}} + 1.0\,\mathcal{L}_{\text{causal-inv}} + 0.5\,\mathcal{L}_{\text{latent-inv}} + 1.0\,D_{\text{JS}} + 0.5\,\mathcal{L}_{\text{DP}} + 0.5\,\mathcal{L}_{\text{EO}},
\end{equation}
where
\begin{itemize}
    \item $\mathcal{L}_{\text{task}}$ is the cross-entropy loss for stress classification;
    \item $\mathcal{L}_{\text{conf}}$ supervises the confounder branch to predict the scar label from $\mathbf{v}_{\text{conf}}$, so that scar information is routed into the confounder subspace;
    \item $\mathcal{L}_{\text{causal-inv}}$ encourages $\mathbf{v}_{\text{causal}}$ to be invariant to the presence of the scar;
    \item $\mathcal{L}_{\text{latent-inv}}$ encourages invariance of the fused representation;
    \item $D_{\text{JS}} = \operatorname{JS}\bigl(f(X) \,\|\, f(X^{S \leftarrow 0})\bigr)$ is the Jensen--Shannon divergence between the predicted distributions for an input and its scar-removed counterfactual; and
    \item $\mathcal{L}_{\text{DP}}$ and $\mathcal{L}_{\text{EO}}$ are differentiable surrogates for the demographic-parity and equalized-odds gaps.
\end{itemize}
% TODO(authors): Give the exact definitions of L_causal-inv, L_latent-inv, L_DP and L_EO
% (distance measure, and which surrogate functions are used for DP and EO).
Because the fairness terms are surrogates, a low training loss does not by itself guarantee the corresponding fairness properties~\cite{yao2024}; all fairness criteria are therefore evaluated directly on held-out data.

Checkpoints are selected on the validation set by
\begin{equation}
\text{score} = \text{Acc} - 0.5\,|\Delta\text{DP}| - 0.5\,\max(\Delta\text{EO}) - 0.2\,\text{CF Gap}.
\end{equation}

\section{Dataset and Experimental Protocol}
% TODO(authors): The markdown describes the physiological input as wrist-worn BVP/EDA
% (Section 3.1) but Appendix A.1 also lists "N=15 WESAD physiology subjects". State
% exactly which dataset supplies the physiological input and, if both are used, how
% they are combined.
The experiments use a pre-registered pilot cohort of eight participants (s1--s8) from UBFC-Phys~\cite{sabour2023}, a multimodal dataset recorded during a social-stress protocol. Each participant contributes three video clips: a rest task (T1), labelled non-stress ($Y = 0$), and speech (T2) and arithmetic (T3) tasks from the Trier Social Stress Test, labelled stress ($Y = 1$). The 24 clips yield 3,344 temporal windows, split into 2,344 training, 504 validation, and 496 test windows. A synthetic brow-line scar is added to the facial frames, with a segmentation mask recorded for each scarred sample; scar-removed counterfactuals are produced from the mask.
% TODO(authors): State whether the train/validation/test split is by participant or by
% window. With eight participants, a window-level split places data from the same
% person in every split and makes test accuracy optimistic.

To test whether the model ignores the scar, the correlation $\rho$ between the scar and the stress label is varied at test time across three regimes:
\begin{itemize}
    \item \textbf{Regime 1 (biased, $\rho = 0.85$):} the scar occurs mostly with stress, as in a biased training set.
    \item \textbf{Regime 2 (neutral, $\rho = 0.50$):} the scar is distributed equally across classes.
    \item \textbf{Regime 3 (inverted, $\rho = 0.15$):} the scar occurs mostly without stress, reversing the training correlation.
\end{itemize}
Each regime is evaluated in both privileged mode, with the ground-truth mask, and autonomous mode, without it.

Four models are compared, each trained with five random seeds (42, 100, 2026, 777, and 888) for 250 epochs with batch size 32 and learning rate $10^{-4}$, all using a MobileNetV3-Small backbone:
\begin{table}[htbp]
\centering
\begin{tabular}{@{}lll@{}}
\toprule
\textbf{Model} & \textbf{Description} & \textbf{Purpose} \\
\midrule
A: Naive ERM & Concatenation fusion, no fairness terms & Tests reliance on the scar \\
B: Camera-off & Physiology-only MLP & Physiological performance reference \\
C: CGF (pruned) & Counterfactual penalty, no Stiefel decomposition & Tests a penalty-only approach \\
D: EQUITAS-RCMF & Stiefel decomposition, LUPI focus, gating & Proposed model \\
\bottomrule
\end{tabular}
\caption{Models compared in the ablation study.}
\label{tab:ablation_models}
\end{table}
% TODO(authors): The markdown calls Model C "CGP Pruned Edge" but its table reports
% "cgf_fair". Use one name consistently (CGF = Causal Gated Fusion).

Training used an RTX~4060 GPU, and the edge benchmark used an x86-64 CPU.
